\chapter{Difference-in-Differences}
\label{cap:did}
% ── index ───────────────────────────────────────────────────────
\index{difference-in-differences|textbf}
\index{DiD|see{difference-in-differences}}
\index{did@\texttt{did()}|textbf}
\index{parallel trends|textbf}
\index{ATT|see{average treatment effect on the treated}}
\index{average treatment effect on the treated|textbf}

% ─────────────────────────────────────────────────────────────────────────────
\section{The model}

Difference-in-Differences (DiD) estimates the causal effect of a treatment
by comparing the time variation of the treated group with that of the control group:

\[
  \hat\delta_{\mathrm{DiD}}
  = \underbrace{(\bar{Y}_{1,\text{post}} - \bar{Y}_{1,\text{pre}})}_{\text{treated variation}}
  - \underbrace{(\bar{Y}_{0,\text{post}} - \bar{Y}_{0,\text{pre}})}_{\text{control variation}}
\]

The $2\times 2$ table summarizes the four groups:

\begin{center}
\begin{tabular}{lcc}
\toprule
 & \textbf{Pre} & \textbf{Post} \\
\midrule
Control ($D=0$) & $\mu_{00}$ & $\mu_{01}$ \\
Treated  ($D=1$) & $\mu_{10}$ & $\mu_{11}$ \\
\bottomrule
\end{tabular}
\end{center}

\[
  \hat\delta = (\mu_{11} - \mu_{10}) - (\mu_{01} - \mu_{00})
\]

% ─────────────────────────────────────────────────────────────────────────────
\section{Regression form}

The DiD estimator is equivalent to the coefficient on the interaction term in:

\[
  y_{it} = \beta_0 + \beta_1\,\text{Post}_t + \beta_2\,\text{Treated}_i
           + \delta\,(\text{Treated}_i \times \text{Post}_t) + \varepsilon_{it}
\]

where $\hat\delta = \hat\beta_3$ is the Average Treatment effect on the Treated (ATT).

% ─────────────────────────────────────────────────────────────────────────────
\section{Parallel trends assumption}

Identification of the ATT rests on:

\[
  E\!\left[Y_{it}^{(0)} - Y_{is}^{(0)} \mid D_i = 1\right]
  = E\!\left[Y_{it}^{(0)} - Y_{is}^{(0)} \mid D_i = 0\right]
  \qquad \forall\; t \neq s
\]

In the absence of treatment, treated and control groups would have the \emph{same}
time trend. The assumption is not testable in the post-treatment period,
but can be checked in pre-treatment periods when more than one
period is available.

\begin{aviso}
  The parallel trends assumption is stronger than it appears. Groups
  with structurally different growth trends (e.g.\ expanding regions
  vs.\ stagnant regions) violate the assumption even if their
  pre-treatment levels are similar.
\end{aviso}

% ─────────────────────────────────────────────────────────────────────────────
\section{DGP Verification}

\subsection*{Exact recovery (no noise)}

The DGP uses 500 individuals $\times$ 2 periods ($n = 1000$). The first 250
are treated, the rest are control. The true effect is $\delta = 3{,}0$:

\[
  y_{it} = 1{,}0 + 0{,}5\cdot\text{Post}_t + 3{,}0\cdot(D_i\times\text{Post}_t)
\]

The four cell means are $\mu_{00}=1{,}0$, $\mu_{01}=1{,}5$,
$\mu_{10}=1{,}0$ and $\mu_{11}=4{,}5$, so that
$\hat\delta = (4{,}5 - 1{,}0) - (1{,}5 - 1{,}0) = 3{,}0$ exactly.

\begin{lstlisting}[caption={Noise-free DiD DGP --- exact recovery}]
seed(42)
input df
id
1
end
for i in 1..=10 { df = append(df, df) }
df |> filter(_n <= 1000)
generate df id      = (_n - 1) % 500 + 1
generate df post    = _n > 500
generate df treated = id <= 250
generate df att     = treated * post
generate df y       = 1.0 + 0.5*post + 3.0*att
did(y ~ treated + post, df)
\end{lstlisting}

\begin{lstlisting}[language={}, basicstyle=\ttfamily\small,
  frame=none, numbers=none, backgroundcolor=\color{codbg},
  xleftmargin=0pt, xrightmargin=0pt, breaklines=false]
================= Difference-in-Differences (2x2 Canonical) ==================
ATT (Effect):                 3.0000 || R-squared:                    1.0000
Std. Error:                   0.0000 || P-value:                      0.0000
t-statistic:                     inf || Observations:                   1000

-------------------------------- Group Means ---------------------------------
Control Group (Pre):      1.0000 | Control Group (Post):     1.5000
Treated Group (Pre):      1.0000 | Treated Group (Post):     4.5000
Parallel Trend Diff:      0.0000 (If > 0, Control grew more than Treated Pre-trend)
==============================================================================
\end{lstlisting}

\subsection*{With idiosyncratic noise}

Adding $\varepsilon \sim \mathcal{N}(0,1)$, DiD remains consistent:

\begin{lstlisting}[caption={DiD DGP with noise}]
generate df e = rnormal()
generate df y = 1.0 + 0.5*post + 3.0*att + e
did(y ~ treated + post, df)
\end{lstlisting}

\begin{lstlisting}[language={}, basicstyle=\ttfamily\small,
  frame=none, numbers=none, backgroundcolor=\color{codbg},
  xleftmargin=0pt, xrightmargin=0pt, breaklines=false]
================= Difference-in-Differences (2x2 Canonical) ==================
ATT (Effect):                 3.0114 || R-squared:                    0.9645
Std. Error:                   0.0358 || P-value:                      0.0000
t-statistic:                 84.0942 || Observations:                   1000

-------------------------------- Group Means ---------------------------------
Control Group (Pre):      1.4842 | Control Group (Post):     2.0099
Treated Group (Pre):      1.4880 | Treated Group (Post):     5.0251
Parallel Trend Diff:      0.0000 (If > 0, Control grew more than Treated Pre-trend)
==============================================================================
\end{lstlisting}

The ATT of $3{,}0114$ is close to the true $3{,}0$; the deviation is sampling
variance. The \emph{Parallel Trend Diff} near zero confirms that the
parallel trends assumption is satisfied by construction.

% ─────────────────────────────────────────────────────────────────────────────
\section{Syntax}

\begin{lstlisting}
did(outcome ~ treated + post, df)
did(outcome ~ treated + post, df, cov=robust)
\end{lstlisting}

The formula requires exactly two variables on the right-hand side: the group
indicator (\lstinline{treated}) and the period indicator (\lstinline{post}).
\hayashi{} computes the interaction internally.

\begin{lstlisting}[caption={Typical DiD workflow}]
load "data.csv" as df

// create indicators
generate df post    = year >= 2020
generate df treated = state == "RS"

let m = did(wage ~ treated + post, df, cov=robust)
display m
\end{lstlisting}

\section{Capturing the result}

\begin{lstlisting}
let m = did(y ~ treated + post, df)
display m
\end{lstlisting}

\begin{nota}
  For event studies with multiple pre- and post-treatment periods,
  the canonical $2\times 2$ DiD may be insufficient. In those cases,
  estimate the regression model with period-by-period interactions and
  verify the absence of pre-treatment effects as a test of
  parallel trends.
\end{nota}

\section{Event study (dynamic DiD)}

\index{eventstudy@\texttt{eventstudy()}}
\index{event study}

The \texttt{eventstudy} command estimates a dynamic DiD with leads and
lags, constructing event-time dummies for each period relative to the
event (excluding the reference period):

\begin{lstlisting}
eventstudy(y ~ event_time + controls, df, ref=-1, min=-5, max=5, cov=HC1)
\end{lstlisting}

\begin{itemize}
  \item \texttt{event\_time}: variable with time relative to the event
    ($\ldots, -2, -1, 0, 1, 2, \ldots$)
  \item \texttt{ref}: reference period omitted (default $-1$)
  \item \texttt{min}, \texttt{max}: lower and upper bounds of event
    times included (default $-5$ and $5$)
  \item Optional controls are included after \texttt{event\_time}
\end{itemize}

The output reports coefficients, standard errors, $t$-statistics, and
$p$-values for each event time. Lead coefficients ($t < 0$) test the
parallel trends assumption: if statistically significant, they indicate
a violation of the identifying assumption.
